% GENERATED FILE. Edit canonical lessons, then run npm run generate:books.
% canonical-source-hash: fnv1a64:bdb222babe8b7a1e
% canonical-lessons: FR-C353-remede, FR-C353-remedier, FR-C353-aggraver, FR-C353-attenuer, FR-C353-compenser

\chapter{Trouver un remède}
\label{ch:fr-trouver-un-remede}
\hlchaptermodality{353}

\begin{chapteropening}
\textbf{By the end of this chapter:} \emph{I can say how to fix a situation: find a remedy, put it right, soften it or make up for it, or say it got worse.}

Complete the last lesson of chapter 353: I can say how to fix a situation: find a remedy, put it right, soften it or make up for it, or say it got worse.
\end{chapteropening}

\section[le remède]{le remède — the remedy, the cure}

\label{lesson:FR-C353-remede}



\begin{quote}
Before the new one: say the French for the asset, then the French for the brake.
\end{quote}

\subsection*{You'll want to know: le remède}
\textbf{le remède} — \textquotedblleft{}the remedy, the cure\textquotedblright{}.

\begin{quote}
Le meilleur remède contre la fatigue, c'est le sommeil.
\end{quote}

\begin{quote}
\emph{The best remedy for tiredness is sleep.}
\end{quote}

\emph{Un remède de grand-mère} — a home remedy.

\begin{cousinweb}
\emph{remède} comes from Latin \emph{remedium} \textquotedblleft{}a cure, a remedy\textquotedblright{} (EtymDB; Lewis \& Short), built from \emph{re-} and \emph{medērī} \textquotedblleft{}to heal\textquotedblright{}.
\end{cousinweb}

\begin{grammarlens}[title={What you've built}]
One more word for telling, explaining and describing at B1.
\end{grammarlens}

\subsection*{Your turn}
\begin{itemize}
\raggedright
  \item \emph{Say it:} \emph{le remède}
  \item \emph{Say it:} \emph{le remède}, once more
  \item \emph{Say it:} a sentence of your own with \emph{le remède}
  \item \emph{Recall:} say the French for the asset, then the French for the brake, then say \emph{le remède} again
\end{itemize}

\begin{tcolorbox}[breakable,colback=teal!4,colframe=teal!35!black,title={Before you move on}]
What does le remède mean? (\textquotedblleft{}The remedy, the cure\textquotedblright{}.) Say it once more.
\end{tcolorbox}
\section[remédier]{remédier — to remedy, to put right}

\label{lesson:FR-C353-remedier}



\begin{quote}
Before the new one: say the French for the brake, then the French for the remedy.
\end{quote}

\subsection*{You'll want to know: remédier}
\textbf{remédier} — \textquotedblleft{}to remedy, to put right\textquotedblright{}.

\begin{quote}
Il faut remédier à ce problème rapidement.
\end{quote}

\begin{quote}
\emph{This problem has to be put right quickly.}
\end{quote}

\emph{Remédier à quelque chose}, with \emph{à}.

\begin{cousinweb}
\emph{remédier} is a learned borrowing from Latin \emph{remediāre} \textquotedblleft{}to cure\textquotedblright{} (EtymDB), from \emph{remedium} \textquotedblleft{}a cure\textquotedblright{}, built on \emph{medērī} \textquotedblleft{}to heal\textquotedblright{} (Lewis \& Short).
\end{cousinweb}

\begin{grammarlens}[title={What you've built}]
One more word for telling, explaining and describing at B1.
\end{grammarlens}

\subsection*{Your turn}
\begin{itemize}
\raggedright
  \item \emph{Say it:} \emph{remédier}
  \item \emph{Say it:} \emph{remédier}, once more
  \item \emph{Say it:} a sentence of your own with \emph{remédier}
  \item \emph{Recall:} say the French for the brake, then the French for the remedy, then say \emph{remédier} again
\end{itemize}

\begin{tcolorbox}[breakable,colback=teal!4,colframe=teal!35!black,title={Before you move on}]
What does remédier mean? (\textquotedblleft{}To remedy, to put right\textquotedblright{}.) Say it once more.
\end{tcolorbox}
\section[aggraver]{aggraver — to make worse, to aggravate}

\label{lesson:FR-C353-aggraver}



\begin{quote}
Before the new one: say the French for the remedy, then the French for to remedy.
\end{quote}

\subsection*{You'll want to know: aggraver}
\textbf{aggraver} — \textquotedblleft{}to make worse, to aggravate\textquotedblright{}.

\begin{quote}
Le froid a aggravé sa toux.
\end{quote}

\begin{quote}
\emph{The cold made his cough worse.}
\end{quote}

\emph{S'aggraver} — to get worse: \emph{la situation s'aggrave}.

\begin{cousinweb}
\emph{aggraver} is a learned borrowing from Latin \emph{aggravāre} \textquotedblleft{}to make heavier\textquotedblright{} (EtymDB), built on \emph{gravis} \textquotedblleft{}heavy\textquotedblright{} (Lewis \& Short).
\end{cousinweb}

\begin{grammarlens}[title={What you've built}]
One more word for telling, explaining and describing at B1.
\end{grammarlens}

\subsection*{Your turn}
\begin{itemize}
\raggedright
  \item \emph{Say it:} \emph{aggraver}
  \item \emph{Say it:} \emph{aggraver}, once more
  \item \emph{Say it:} a sentence of your own with \emph{aggraver}
  \item \emph{Recall:} say the French for the remedy, then the French for to remedy, then say \emph{aggraver} again
\end{itemize}

\begin{tcolorbox}[breakable,colback=teal!4,colframe=teal!35!black,title={Before you move on}]
What does aggraver mean? (\textquotedblleft{}To make worse, to aggravate\textquotedblright{}.) Say it once more.
\end{tcolorbox}
\section[atténuer]{atténuer — to lessen, to soften}

\label{lesson:FR-C353-attenuer}



\begin{quote}
Before the new one: say the French for to remedy, then the French for to make worse.
\end{quote}

\subsection*{You'll want to know: atténuer}
\textbf{atténuer} — \textquotedblleft{}to lessen, to soften\textquotedblright{}.

\begin{quote}
Ce médicament atténue la douleur.
\end{quote}

\begin{quote}
\emph{This medicine eases the pain.}
\end{quote}

\emph{S'atténuer} — to fade, to ease.

\begin{cousinweb}
\emph{atténuer} is a learned borrowing from Latin \emph{attenuāre} \textquotedblleft{}to make thin, to weaken\textquotedblright{} (Lewis \& Short), built on \emph{tenuis} \textquotedblleft{}thin\textquotedblright{}.
\end{cousinweb}

\begin{grammarlens}[title={What you've built}]
One more word for telling, explaining and describing at B1.
\end{grammarlens}

\subsection*{Your turn}
\begin{itemize}
\raggedright
  \item \emph{Say it:} \emph{atténuer}
  \item \emph{Say it:} \emph{atténuer}, once more
  \item \emph{Say it:} a sentence of your own with \emph{atténuer}
  \item \emph{Recall:} say the French for to remedy, then the French for to make worse, then say \emph{atténuer} again
\end{itemize}

\begin{tcolorbox}[breakable,colback=teal!4,colframe=teal!35!black,title={Before you move on}]
What does atténuer mean? (\textquotedblleft{}To lessen, to soften\textquotedblright{}.) Say it once more.
\end{tcolorbox}
\section[compenser]{compenser — to make up for, to compensate}

\label{lesson:FR-C353-compenser}



\begin{quote}
Before the new one: say the French for to make worse, then the French for to lessen.
\end{quote}

\subsection*{You'll want to know: compenser}
\textbf{compenser} — \textquotedblleft{}to make up for, to compensate\textquotedblright{}.

\begin{quote}
Il travaille le samedi pour compenser.
\end{quote}

\begin{quote}
\emph{He works on Saturdays to make up for it.}
\end{quote}

\emph{Compenser le temps perdu} — to make up for lost time.

\begin{cousinweb}
\emph{compenser} is a learned borrowing from Latin \emph{compēnsāre} \textquotedblleft{}to weigh one thing against another\textquotedblright{} (EtymDB; Lewis \& Short), built on \emph{pēnsāre} \textquotedblleft{}to weigh\textquotedblright{}.
\end{cousinweb}

\begin{grammarlens}[title={What you've built}]
One more word for telling, explaining and describing at B1.
\end{grammarlens}

\subsection*{Your turn}
\begin{itemize}
\raggedright
  \item \emph{Say it:} \emph{compenser}
  \item \emph{Say it:} \emph{compenser}, once more
  \item \emph{Say it:} a sentence of your own with \emph{compenser}
  \item \emph{Recall:} say the French for to make worse, then the French for to lessen, then say \emph{compenser} again
\end{itemize}

\begin{tcolorbox}[breakable,colback=teal!4,colframe=teal!35!black,title={Before you move on}]
What does compenser mean? (\textquotedblleft{}To make up for, to compensate\textquotedblright{}.) Say it once more.
\end{tcolorbox}
